% Post-CP23 candidate. This is a fixed-source geometric certificate, not a uniform image theorem.
\subsection{Twenty nodal five-self-isogenies of a supersingular Jacobian}
\label{subsec:five-theta-self-isogenies-candidate}
\status{Proved and independently checked.}

Let $k=\overline{\mathbf F}_2$, let $C$ be the smooth projective curve
\[
                         y^2+y=x^5,
\]
let $W$ be its unique point at infinity, and embed $C$ as the symmetric
principal theta divisor of $J=\operatorname{Pic}^0(C)$ by $P\mapsto[P-W]$.
Write $\mathbf F_{16}=\mathbf F_2[t]/(t^4+t+1)$, with field elements
encoded by their four binary coefficients. The Cartier matrix on
$\langle dx,x\,dx\rangle$ is nilpotent, since $\mathcal C(dx)=0$
and $\mathcal C(x\,dx)=dx$. Thus $f(J)=0$ and $J[2](k)=0$.

\begin{lemma}[Explicit fifth-order divisor functions and pairings]
\label{lem:five-theta-explicit-Weil-candidate}
\textup{[Proved; independent audits passed.]}
The curve has thirty-three $\mathbf F_{16}$-points. Its thirty-two
finite points lie in $J[5]$, and they generate the entire four-dimensional
symplectic $\mathbf F_5$-space $J[5](k)$. In particular every geometric
fifth-torsion point of $J$ is $\mathbf F_{16}$-rational.
For distinct finite points $P=(a,b)$ and $Q$, put
\[
             f_P=y+a^8x^2+a^4x+b^4.
\]
Then
\[
 \operatorname{div}(f_P)=5P-5W,\qquad
 e_5(P-W,Q-W)=\frac{f_P(Q)}{f_Q(P)}.
\]
The inverse convention for $e_5$ gives the same isotropic subspaces.
\end{lemma}
\begin{proof}
For $a\in\mathbf F_{16}$, $a^5\in\mathbf F_4$, whose absolute trace
from $\mathbf F_{16}$ is zero. Hence $b^2+b=a^5$ has two solutions
in $\mathbf F_{16}$, giving the thirty-two finite points.
Since $b^4=b+a^{10}+a^5$, direct substitution gives
\[
                      f_P^2+f_P=(x+a)^5.
\]
At $P$ the function vanishes, at the other point above $a$ it is one,
and its only pole is $W$, of order five. This proves its divisor.
These are the explicit functions underlying
\cite[Example~1.2(ii), pp.~30--31]{BoxallGrantLeprevost2001}.

Here is the common-pole normalization in the pairing formula.
The parameter $z=x^2/y$ is a uniformizer at $W$ and
$z^5 f_P\to1$, independently of $P$. In the divisor definition of
the Weil pairing, move $Q-W$ to a divisor disjoint from $P-W$.
For completeness choose a rational function $h$ having a simple zero
at $W$ and no zero or pole at $P,Q$; weak approximation supplies one.
If $h\sim d z$ and $f_P\sim c_Pz^{-5}$ at $W$, the tame-symbol
product formula gives
\[
           f_P(\operatorname{div}(h)-W)
                   =\frac{h(P)^5}{c_Pd^5}.
\]
Indeed the symbols at $P,W$ are respectively $h(P)^{-5}$ and
$c_Pd^5$, and all other symbols give the evaluation on the left.
Using the disjoint divisor $Q-W+\operatorname{div}(h)$ and its
function $f_Qh^5$ therefore gives
$f_P(Q)c_Q/(f_Q(P)c_P)$. Here $c_P=c_Q=1$; the usual sign is also
one in characteristic two. This proves the displayed formula.

The retained exact program enumerates the thirty-three points,
evaluates these functions, and expresses all pairings as powers of
$\zeta=t^3$. Four finite points, numbered $1,3,5,9$ in its displayed
point list, have pairing matrix
\[
 J_5=\begin{pmatrix}
 0&0&4&3\\0&0&4&2\\1&1&0&3\\2&3&2&0
 \end{pmatrix}\quad\text{over }\mathbf F_5.
\]
Its determinant is nonzero, so their classes are independent.
The group $J[5](k)$ has dimension four; these rational classes are
therefore a basis. The program reconstructs the coordinates of
every curve point from its four pairings and verifies all pairings
again in those coordinates. In particular this computation does
not assume a Jacobian group-law formula.
\end{proof}

\begin{lemma}[A complete geometric theta-packet census]
\label{lem:five-theta-geometric-census-candidate}
\textup{[Proved; independent audits passed.]}
Among the $156$ maximal isotropic subgroups $K\subset J[5]$, exactly
twenty give a geometrically integral nodal image of $C$ in $J/K$.
Each such image has twenty-four ordinary nodes. For the other
subgroups, the largest rational packet has size four for eighty
subgroups and size five for fifty-six subgroups.
\end{lemma}
\begin{proof}
For a nonzero $g\in J(k)$, $C\cap(C-g)$ has length two: theta has
self-intersection two and no nontrivial translation stabilizer.
If $P,P+g\in C$, symmetry also gives $-P-g,-P\in C$.
These describe the two intersection points, with their possible
coincidence. Suppose a geometric packet contains three distinct
points $P,P+g,P+h$, where $g,h\in J[5]$ are distinct and nonzero.
Both intersection schemes are defined over $\mathbf F_{16}$.
If $P$ were not rational, its Frobenius image would have to be
the other point of each length-two scheme. Thus
\[
             \operatorname{Frob}_{16}(P)=-P-g=-P-h,
\]
which contradicts $g\ne h$. Consequently every geometric packet
of size at least three is already in the rational point list.

The tangent directions of this theta embedding are the fibres of
the canonical degree-two map of $C$. Hence distinct identified
points can have equal tangents only when they are $P$ and $-P$.
If $2P\in K\subset J[5]$, then $[5]P\in J[2](k)=0$; therefore
$P\in J[5](\mathbf F_{16})$. Moreover multiplication by two is
invertible on $K$, so $P\in K$. The packet containing $W$ then
contains the three distinct points $W,P,-P$. Thus a rational
census with every packet of size at most two also excludes every
geometric tangency.

The program enumerates the $156$ projective lines in
$\mathbf F_5^4$, and the $156=(5+1)(5^2+1)$ isotropic planes by
spanning every pair of independent orthogonal lines. For each
plane it partitions the thirty-three curve vectors into cosets.
It retains every plane basis and every packet, and finds
\[
 \begin{array}{c|ccc}
 \text{largest packet}&2&4&5\\ \hline
 \text{number of planes}&20&80&56.
 \end{array}
\]
The preceding geometric argument proves that the twenty first
planes have only transverse double identifications, and that the
other planes cannot have nodal images.

The quotient is etale. Its restriction to $C$ is birational onto
its image: a positive-degree generic identification would make
$C$ invariant under a nonzero element of $K$, contrary to its
principal polarization. Thus the image is integral with
normalization $C$, and the identified transverse pairs give exactly
ordinary nodes. Its numerical class on the principally polarized
quotient is $5\Theta'$, so its arithmetic genus is $26$ and the
node number is $26-2=24$.
\end{proof}

\begin{proposition}[Explicit nodal polarized five-self-isogenies]
\label{prop:five-theta-twenty-self-isogenies-candidate}
\textup{[Proved; independent audits passed.]}
Let $\alpha,\beta\in\mathbf F_{16}$ satisfy
\[
        \alpha^5\in\mathbf F_4\setminus\mathbf F_2,
        \qquad \beta^2+\beta=\alpha^5.
\]
The automorphism
\[
 g_{\alpha,\beta}(x,y)
   =(x+\alpha,\ y+\alpha^8x^2+\alpha^4x+\beta)
\]
induces an automorphism $g$ of the principal Jacobian with
$g^2=[-1]$. The isogeny $\phi=[2]+g$ has degree twenty-five and
$\phi^*\Theta\equiv5\Theta$. Its kernel is one of the twenty good
planes of the preceding lemma. As $(\alpha,\beta)$ varies, these
twenty kernels are distinct and exhaust the good planes.
Consequently every such $\phi$ gives a nodal theta image on this
same fixed Jacobian, defined over $\mathbf F_{16}$.
\end{proposition}
\begin{proof}
Substitution in the curve equation verifies the automorphism.
Its square fixes $x$ and adds
$\alpha^{10}+\alpha^5=1$ to $y$, so it is the hyperelliptic
involution. It fixes $W$ and preserves the principal polarization;
its Rosati adjoint is $g^{-1}=-g$. Therefore
\[
                 \phi^\dagger\phi=(2-g)(2+g)=[5].
\]
It is a separable polarized isogeny of degree twenty-five, and its
kernel is maximal isotropic in $J[5]$.

There are ten possible $\alpha$ and two possible $\beta$ for each.
The second retained program computes their actions on the four
explicit fifth-torsion basis points of the first lemma. It verifies
the entire action on all thirty-three points, $g^2=-I$, and
$g^{\mathsf t}J_5g=J_5$. It then enumerates
$\ker(2I+g)\subset\mathbf F_5^4$ and verifies that the twenty
resulting order-twenty-five subgroups are distinct and exactly
the twenty good planes of the complete census.

The norm divisor is symmetric. Its line bundle differs from
$\mathcal O_J(5\Theta)$ by a symmetric element of $\operatorname{Pic}^0(J)$,
which is a geometric two-torsion point and hence zero. Thus these
are actual nodal divisors in the indicated complete system,
not merely numerical classes.
\end{proof}

\begin{remark}[Certificate and remaining monodromy scope]
\label{rem:five-theta-self-isogeny-scope-candidate}
The executed programs are
\path{code/monodromy/nonordinary_five_theta_packet_probe.py} and
\path{code/monodromy/nonordinary_five_theta_self_isogenies.py}, with complete
JSON outputs of the same stems in \path{agents/monodromy/outputs/}. The first passed
at 20:08:39 UTC in $0.31384029000037117$ seconds; an independent
output-preserving rerun passed in $0.2980953239966766$ seconds.
The second passed at 20:24:45 UTC in $0.051931495996541344$ seconds.
The parent completed the full mathematical and source audit at 20:31:15 UTC and independently reran both complete programs, preserving the author outputs, in $0.334064817$ and $0.058760470$ seconds. The independent reduction/descent audit also checked the Weil functions and finite census; no broader geometric scope is attributed to that bounded audit.
This establishes good degree-five sources and targets only for the
specified Jacobian. It does not establish a suitable five-neighbour
for every fixed nonordinary surface, nor the required degree-five
node permutation and gluing image. The order-three eigenmap
argument cannot be imported for a fifth-torsion translation,
which need not preserve $\mathcal O(3\Theta)$.
\end{remark}
